\section{Full zero-head S from one-arm C, using a double-defect dichotomy}
\begin{theorem}
Let $Z,Y$ be actual normalized increasing arms, each with at least two
non-one parts. Assume at least one arm has full property C. Then the
correctly conditioned coupled S resource holds at EVERY zero-head
crossing $k=0$:
\[
                  u_0+v_n\le2(u_1+v_{n+1}).              \tag{1}
\]
All original S hypotheses, including strict weighted fall and
$2p_n\ge p_{n-1}$, are retained. No OLD A2 assumption is needed
at this particular head, since $u_0=1$, $u_1=p_1-1$ and $p_1\ge6$.
Consequently (1) is unconditional whenever either old arm has at most
five non-one parts, by the completed C5 theorem; the other arm is
unrestricted. General original coefficient unimodality is not asserted.
\end{theorem}
\begin{proof}
Write $A=F_Z$, $B=F_Y$, $P=AB$, $D=\operatorname{wt}Z+
\operatorname{wt}Y$, $p(z)=z^DP(1/z)$, and $j=D-n$.
Strict weighted fall at $k=0$, together with the nonnegative tail,
gives the integer bound
\[
                    e:=2P_j-P_{j+1}\in\{0,1\}.          \tag{2}
\]
If $j<\kappa_Y$, the local opposite-universality theorem applies,
including its arbitrary first-defect boundary at $j+1$. If
$j<\kappa_Z$, exchange the arms in the same coefficient argument.
That argument uses only the coefficient conditions, so the physical
orientation of the append need not be exchanged. These cases prove (1).

It remains that $j\ge\max(\kappa_Z,\kappa_Y)$.
For any actual nonempty arm $Y$, the forward deficit $E_Y$ has positive
interval support $[\kappa_Y,\operatorname{wt}Y+1]$, with all positive
coefficients integral and terminal coefficient two. Every coefficient
of $A$ on $[0,\operatorname{wt}Z]$ is at least one. Therefore
\[
          (E_YA)_{j+1}\ge2,\qquad(E_ZB)_{j+1}\ge2.       \tag{3}
\]
Here $j+1$ is past each first support point and no later than the
product's final degree. To verify the elementary convolution bound:
if the deficit support has one point, that coefficient is two and
every product coefficient is at least two. Otherwise every interior
convolution rank past its first point has at least two positive
integer summands. At the final rank its sole summand uses the terminal
deficit coefficient two. Thus all boundary cases of (3) are included.

The two exact identities
\[
 (2z-1)P=E_YA-(1-z)A=E_ZB-(1-z)B
\]
now give, using (2)--(3),
\[
                 A_{j+1}-A_j\ge1,\qquad B_{j+1}-B_j\ge1. \tag{4}
\]
Both ordinary factors are strictly rising at $j+1$; their positive
differences at $j$ are also positive by initial interval support.
Also $j\ge\max(\kappa_Z,\kappa_Y)\ge3$.
Apply the C-to-S rising-window theorem, choosing whichever arm has C.
It gives $R_{n+1}\le R_n$, where $R_i=2p_i-p_{i-1}$.
The head budget $2u_1-u_0=2p_1-3>0$ then proves (1).
\end{proof}

This dichotomy makes the auxiliary C advance contribute to the original
weighted-difference induction. For a possible properly conditioned S
failure with both old arms having at least two non-one parts, zero head
now requires BOTH old arms to lie outside the C5 range. At nonzero head,
the C-to-S bridge still closes the two-rising-factor region whenever
one arm has C and the OLD A2 inductive premise holds. Neither statement
raises the original conjecture's previously proved least-counterexample
non-one count: the remaining original obstruction may have nonzero head.

The three reports in \texttt{CToCoupledSRisingWindow.lean} include
the complete scalar double-defect implication in (3)--(4), in addition
to the monotone-excess and S-sum lemmas. The local first-defect proof
has four scoped reports in \texttt{UntruncatedOppositeZeroHead.lean}.
The positive convolution bound and exact actual word assembly are
the complete written argument above, not an end-to-end formalization.
General S at nonzero head outside the rising region, unrestricted C
beyond five, A2 and the original unamplified conjecture remain open.
